% ============================================================================
% Seccion 8: Resultados y Demostracion
% TFP - Rubrica: 2 puntos
% ============================================================================
\section{Resultados y demostraci\'on}
\label{sec:results}

\subsection{Estado del prototipo}

El c\'odigo del proyecto se organiza en la organizaci\'on de GitHub
\texttt{spark-match}, con repositorios separados por dominio (datos, backend,
agente, frontend, infraestructura, DevOps y art\'iculo). A la fecha est\'an
implementados:

\begin{itemize}[noitemsep]
  \item el \emph{pipeline de datos} que produce \texttt{features.csv} (6\,208
        combinaciones carrera--instituci\'on, 554 carreras \'unicas), con las
        variables de ingreso, costo, admisi\'on y duraci\'on ya normalizadas;
  \item el \emph{etiquetado RIASEC} de las carreras mediante el LLM en Bedrock;
  \item el \emph{agente conversacional} (Deep Agent) con cuatro herramientas
        operativas: evaluaci\'on de perfil RIASEC (\texttt{evaluate\_riasec\_}
        \texttt{profile}), b\'usqueda de carreras (\texttt{search\_careers}),
        c\'alculo de afinidad vocacional (\texttt{calculate\_affinity}) y
        b\'usqueda web de respaldo (\texttt{web\_search}), trazadas con
        LangSmith y con un framework de evaluaci\'on;
  \item los \emph{scaffolds} del backend serverless (TypeScript, contexto
        \texttt{identity}: autenticaci\'on, usuarios y auditor\'ia) y del
        frontend (Angular 21);
  \item la \emph{infraestructura} en Terraform con ambientes \texttt{dev} y
        \texttt{prod}.
\end{itemize}

\paragraph{Alcance real del motor de scoring (importante).} De los cinco
criterios dise\~nados (afinidad, ingreso, costo, admisi\'on, duraci\'on), a la
fecha \textbf{solo la afinidad vocacional (RIASEC) est\'a integrada} en la
herramienta \texttt{calculate\_affinity} del agente, y \'esta opera sobre un
\textbf{cat\'alogo piloto de 19 carreras} definidas manualmente en
\texttt{data/careers/} (con nombre, campo y perfil RIASEC, sin universidad,
ingreso, costo ni admisi\'on) --- \textbf{no} sobre las 6\,208 combinaciones de
\texttt{features.csv}. El backend serverless, por su parte, solo implementa el
contexto \texttt{identity}; no existe a\'un un contexto de \emph{matching} ni
una carga de \texttt{features.csv} en RDS PostgreSQL. En consecuencia, el
ranking multicriterio y el ejemplo num\'erico de la
Secci\'on~\ref{sec:implementation} describen el dise\~no objetivo del sistema,
no un resultado ya medido end-to-end. El trabajo pendiente para cerrar esta
brecha --- integrar ingreso, costo, admisi\'on y duraci\'on al agente, y
conectar su cat\'alogo con el dataset completo --- se detalla en
\texttt{DOCS/issues/issues-v1/}.

\subsection{Ejemplo de salida (caso can\'onico)}

\textbf{Caso:} estudiante de Cusco, presupuesto m\'aximo S/600/mes, prioridad en
salario y cercan\'ia geogr\'afica.

\textbf{Mensaje del usuario:} \textit{``Vivo en Cusco y mi presupuesto m\'aximo es
600 soles al mes. Me interesa algo con buenos sueldos pero tambi\'en quiero
quedarme cerca de casa''.}

\textbf{Preferencias extra\'idas por el LLM:}
\begin{lstlisting}[style=jsonstyle, caption={JSON de preferencias extraidas}, label={lst:json-prefs}]
{
  "region": "Cusco",
  "presupuesto_max": 600,
  "peso_afinidad": 0.15,
  "peso_ingreso": 0.45,
  "peso_costo": 0.30,
  "peso_admision": 0.05,
  "peso_duracion": 0.05
}
\end{lstlisting}

El motor de scoring aplica los filtros duros (regi\'on Cusco, presupuesto) y
ordena las combinaciones carrera--universidad por su puntaje, generando un ranking
Top-N. Finalmente, el agente redacta una justificaci\'on en lenguaje natural que
cita los datos que sustentan cada recomendaci\'on (ingreso promedio, costo,
afinidad).

\begin{figure}[H]
    \centering
    % \includegraphics[width=0.95\linewidth]{chat-screenshot.png}
    \caption{Captura de la interfaz conversacional de Spark Match (por incluir).}
    \label{fig:chat}
\end{figure}
% TODO (equipo): reemplazar por capturas reales de la app corriendo, el ranking
% Top-N de un caso real y el enlace a la demo/video.

\subsection{Evidencia funcional}

La demostraci\'on final mostrar\'a el flujo de extremo a extremo: el estudiante
conversa con el orientador, el agente construye su perfil y prioridades, el motor
de scoring genera el ranking Top-N sobre datos reales y el agente entrega las
recomendaciones explicadas. La evidencia incluir\'a:

\begin{itemize}[noitemsep]
  \item \textbf{Repositorios:} \url{https://github.com/spark-match}.
  \item \textbf{Demo / video:} enlace por definir tras el despliegue.
  \item \textbf{Capturas:} conversaci\'on real con el agente (Figura~\ref{fig:chat}),
        ranking Top-N y trazas de observabilidad en LangSmith.
\end{itemize}
